\clearpage
\begingroup
\fancyhead[L]{\small NACHLASS UND BRIEFWECHSEL}
\renewcommand{\headrulewidth}{0pt}
\par\medskip
\Needspace{8\baselineskip}
\pdfbookmark[2]{[2] Nicolai an Gauss, 19. März 1814}{pallas-VIII-letter2}
\begin{center}
[2.]\par\medskip
\textsc{Nicolai} an \textsc{Gauss}, Seeberg, 19. März 1814.
\end{center}
\par\noindent\hspace*{\parindent}\{.\ .\ .\ .\ .\quad Seit kurzem habe ich nun auch wieder die Saturnsstörungen der
Pallas vorgenommen, und ich mache mir das Vergnügen, Ihnen heute drei
vollendete Elemente mitzutheilen. Sie werden Sich vielleicht wundern über
die ziemliche Anzahl der Gleichungen, worüber ich anfangs selbst erstaunte;
allein das meiste ist Quarkzeug, was die Resultate wohl wenig ändern wird.
Indess habe ich \textls[100]{keine einzige} Gleichung vernachlässigt, die bei \(\log a\) mehr
als eine Einheit, und bei den andern Elementen mehr als \(0\overset{\prime\prime}{,}1\) beträgt, so wie
Sie es beim Jupiter gemacht haben. Auch kann ich die Richtigkeit einer
\textls[100]{jeden} Gleichung verbürgen, da ich sie alle zuerst mit grosser Sorgfalt und
Emsigkeit berechnet, und dann noch einmal flüchtig nachgesehen habe. Ich
glaubte dies thun zu müssen, da man, die Interpolationen etwa ausgenommen,
bei der ganzen Rechnung keine Controlle hat.\quad .\ .\ .\ .\ .\quad Noch muss ich bemerken,
dass ich für die Saturnsmasse diejenige angenommen habe, welche
\textsc{Laplace} nach neuern Untersuchungen in der neuesten 4\textsuperscript{ten} Ausgabe seiner
Exposition du Système du monde, welche vor ein paar Wochen hier ankam, gegeben
hat, nemlich \(\frac{1}{3512{,}08}\mathpunct{\raisebox{.45ex}{.}}\) --- Die Rechnung nach dieser andern Methode
\SourceLeaf{586}
scheint mir übrigens bei weitem einfacher und bequemer, als die nach der
erstern; auch ist man ohne Zweifel weit wenigern Irrungen ausgesetzt; nur
vor dem zweiten Theil der Epoche habe ich etwas Furcht. --- Die Säcularänderungen
für Neigung und Knoten, so wie auch die grössern Gleichungen,
stimmen recht gut mit meiner ersten Rechnung. --- Ich sollte nicht denken,
dass des Periheliums so wie auch des 2\textsuperscript{ten} Theils der Epoche wegen die angenommene
Anzahl von zwölf zum Grunde gelegten Pallasörtern zu klein sein
wird; aber wenn ich die Anzahl der Gleichungen für diese beiden Elemente
nach meiner frühern Rechnung mit der für die andern Elemente nach derselben
Rechnung vergleiche, und nun die Menge der Gleichungen für die
letztern erwäge, welche die neue Rechnung gegeben hat, so fürchte ich es
beinahe. Dieserwegen wünschte ich jetzt, dass ich wenigstens 18 Pallasörter
zum Grunde gelegt hätte. Die Folge der Rechnung wird dies alles weiter
ausweisen. Die numerische Berechnung aller Gleichungen für die 9 Oppositionen
wird übrigens auch noch ein kleines Stück Arbeit sein.\quad .\ .\ .\ .\ .
\par\medskip
\Needspace{28\baselineskip}
\begin{center}
\(\operatorname{Log}a\) und erster Theil der Epoche\HistoricalDisplayNoteMark.
\end{center}
\HistoricalDisplayNoteText{Der Ausdruck links gibt die Störungen von \(\log a\) in Einheiten der
7. Decimale; die rechts stehenden Coefficienten geben, mit dem
entsprechenden Cosinus verbunden, die Störungen \(\delta\int n\,\mathrm{d}t\).]}
\begingroup
\small
\newcommand{\VIIIpallas}{\text{\AstroSymbols ⚴}}
\newcommand{\VIIIsaturn}{\text{\AstroSymbols ♄}}
\setlength{\arraycolsep}{2pt}
\renewcommand{\arraystretch}{1.25}
\[
\begin{array}{r@{\,}l@{\qquad\quad}r}
5&\sin(\VIIIpallas+79^0\ 15')&0\overset{\prime\prime}{,}39\\
+7&\sin(2\VIIIpallas+93\ 48)&0{,}24\\
+2&\sin(3\VIIIpallas+341\ 19)&0{,}04\\
+2&\sin(\VIIIpallas+\VIIIsaturn+87\ 31)&0{,}15\\
+2&\sin(\VIIIpallas-\VIIIsaturn+92\ 48)&0{,}20\\
+2&\sin(2\VIIIpallas-\VIIIsaturn+77\ 14)&0{,}07\\
+2&\sin(3\VIIIpallas-\VIIIsaturn+88\ 13)&0{,}05\\
+1&\sin(\VIIIpallas+2\VIIIsaturn+140\ 23)&0{,}06\\
+38&\sin(\VIIIpallas-2\VIIIsaturn+34\ 30)&3{,}92\\
+39&\sin(2\VIIIpallas-2\VIIIsaturn+88\ 29)&1{,}65\\
+9&\sin(3\VIIIpallas-2\VIIIsaturn+331\ 1)&0{,}23\\
+2&\sin(4\VIIIpallas-2\VIIIsaturn+190\ 0)&0{,}03\\
+10&\sin(\VIIIpallas-3\VIIIsaturn+84\ 44)&1{,}37\\
+6&\sin(2\VIIIpallas-3\VIIIsaturn+84\ 30)&0{,}27
\end{array}
\]
\SourceLeaf{587}
\[
\begin{array}{r@{\,}l@{\qquad\quad}r}
+9&\sin(3\VIIIpallas-3\VIIIsaturn+81^0\ 4')&0\overset{\prime\prime}{,}24\\
+3&\sin(4\VIIIpallas-3\VIIIsaturn+325\ 50)&0{,}06\\
+1&\sin(\VIIIpallas-4\VIIIsaturn+113\ 8)&0{,}28\\
+2&\sin(2\VIIIpallas-4\VIIIsaturn+81\ 57)&0{,}09\\
+2&\sin(3\VIIIpallas-4\VIIIsaturn+82\ 45)&0{,}05\\
+2&\sin(4\VIIIpallas-4\VIIIsaturn+72\ 40)&0{,}04.
\end{array}
\]
\Needspace{10\baselineskip}
\begin{center}
Neigung der Bahn.\par
Jährliche Änderung \(=+0\overset{\prime\prime}{,}16950\)
\end{center}
\[
\begin{array}{r@{\,}l@{\qquad}r@{\,}l}
0\overset{\prime\prime}{,}42&\sin(\VIIIpallas+42^0\ 58')&+0\overset{\prime\prime}{,}11&\sin(\VIIIpallas+2\VIIIsaturn+129^0\ 41')\\
+0{,}24&\sin(2\VIIIpallas+98\ 47)&+0{,}56&\sin(\VIIIpallas-2\VIIIsaturn+243\ 21)\\
+1{,}02&\sin(\VIIIsaturn+37\ 58)&+0{,}29&\sin(2\VIIIpallas-2\VIIIsaturn+275\ 8)\\
+0{,}18&\sin(\VIIIpallas+\VIIIsaturn+89\ 3)&+0{,}23&\sin(3\VIIIsaturn+83\ 5)\\
+0{,}11&\sin(\VIIIpallas-\VIIIsaturn+355\ 47)&+0{,}52&\sin(\VIIIpallas-3\VIIIsaturn+260\ 41)\\
+2{,}27&\sin(2\VIIIsaturn+105\ 16)&+0{,}13&\sin(\VIIIpallas-4\VIIIsaturn+273\ 9).
\end{array}
\]
\Needspace{12\baselineskip}
\begin{center}
Länge des aufsteigenden Knotens.\par
Jährliche Änderung \(=-1\overset{\prime\prime}{,}47934\)
\end{center}
\[
\begin{array}{r@{\,}l@{\qquad}r@{\,}l}
0\overset{\prime\prime}{,}89&\sin(\VIIIpallas+281^0\ 36')&+1\overset{\prime\prime}{,}22&\sin(\VIIIpallas-2\VIIIsaturn+111^0\ 22')\\
+0{,}39&\sin(2\VIIIpallas+11\ 13)&+0{,}59&\sin(2\VIIIpallas-2\VIIIsaturn+184\ 51)\\
+1{,}79&\sin(\VIIIsaturn+246\ 49)&+0{,}31&\sin(3\VIIIsaturn+327\ 21)\\
+0{,}17&\sin(\VIIIpallas+\VIIIsaturn+350\ 51)&+0{,}36&\sin(\VIIIpallas-3\VIIIsaturn+3\ 31)\\
+0{,}75&\sin(\VIIIpallas-\VIIIsaturn+186\ 33)&+0{,}14&\sin(2\VIIIpallas-3\VIIIsaturn+143\ 10)\\
+0{,}13&\sin(2\VIIIpallas-\VIIIsaturn+301\ 7)&+0{,}11&\sin(3\VIIIpallas-3\VIIIsaturn+180\ 41)\\
+3{,}27&\sin(2\VIIIsaturn+2\ 51)&+0{,}13&\sin(\VIIIpallas-4\VIIIsaturn+25\ 2).\}\\
+0{,}15&\sin(\VIIIpallas+2\VIIIsaturn+37\ 19)&&
\end{array}
\]
\endgroup
\clearpage
\endgroup
